% Generated draft: requires mathematical and source review.
\documentclass[11pt]{article}
\usepackage[margin=25mm]{geometry}
\usepackage{fontspec}
\setmainfont{Latin Modern Roman}
\usepackage{amsmath,amssymb,mathtools,unicode-math}
\setmathfont{Latin Modern Math}
\usepackage{enumitem,needspace,xurl,hyperref}
\setlength{\parindent}{0pt}
\setlength{\parskip}{5pt}
\setlength{\emergencystretch}{3em}
\newcommand{\N}{\mathbb{N}}
\newcommand{\Z}{\mathbb{Z}}
\newcommand{\Q}{\mathbb{Q}}
\newcommand{\R}{\mathbb{R}}
\newcommand{\C}{\mathbb{C}}
\newcommand{\F}{\mathbb{F}}
\newcommand{\A}{\mathbb{A}}
\newcommand{\PP}{\mathbb{P}}
\newcommand{\E}{\mathbb{E}}
\newcommand{\eps}{\varepsilon}
\newcommand{\dd}{\,\mathrm{d}}
\begin{document}
\section*{UnsolvedMath ID: 800038}
\textbf{Source problem code:} catalogue entry
\subsection{Definitions and mathematical statement}
A semigroup is a set equipped with an associative binary operation. A semigroup variety is a class of semigroups defined by identities. The supplied problem uses the term \emph{inclusive variety} without defining it; consequently, the following formulation is conditional on recovering that term's intended meaning from the source.

For an inclusive semigroup variety \(\mathcal V\), write
\[
\mathcal I(\mathcal V)=\{\mathcal W\subseteq\mathcal V: \mathcal W\text{ is an inclusive semigroup variety}\},
\]
ordered by inclusion. The source calls this ordered collection the lattice of inclusive subvarieties; its lattice structure cannot be specified without the missing definition of inclusiveness.

\textbf{Problem (G. Mashevitzky).} Determine all inclusive varieties \(\mathcal V\) for which
\[
|\mathcal I(\mathcal V)|<\aleph_0
\qquad\text{or}\qquad
|\mathcal I(\mathcal V)|=\aleph_0.
\]
Thus the requested classification includes both finite and countably infinite lattices of inclusive subvarieties. The source reports a partial solution in reference [M2], whose bibliographic entry and mathematical content are absent from the supplied record.
\subsection{Short English statement}
Classify the inclusive varieties whose lattices of inclusive subvarieties are finite or countably infinite. The meaning of ``inclusive'' must be recovered before this becomes a fully specified mathematical problem. The supplied record labels the problem open and reports a partial solution in [M2]; neither assessment is independently verified here.
\subsection{Sources}
\begin{itemize}
\item Local UnsolvedMath registry, permanent ID \(800038\), source code \texttt{catalogue entry}, title \emph{The Lyapin's Notebook --- Problem 40}. The supplied registry records status ``open'' and proposed year 2026; these are source metadata rather than independently verified historical or mathematical conclusions.
\item Bershadsky, Kublanovsky, and Mashevitzky, \emph{The Lyapin's Notebook: Unsolved Problems in Semigroup Theory} (2026), Problem 40. Supplied source URL: \url{https://arxiv.org/abs/2604.04763}. The record reports access on 2026-10-16 and extraction by \texttt{arxiv\_latex}; these are supplied provenance details, not a claim of access or verification here.
\item The problem statement attributes the question to G. Mashevitzky. The registry's \texttt{proposed\_by} field lists Bershadsky, Kublanovsky, and Mashevitzky.
\item Reference [M2] is cited in the supplied statement as containing a partial solution. Its bibliographic entry is not supplied, so its identity, scope, and conclusions remain unresolved.
\end{itemize}
\end{document}
